\chapter{Many qubits and entanglement}\label{ch:entanglement}

So far we have seen that a lot of useful things can be done already with a single qubit: random numbers and quantum cryptography are part of the
larger field of quantum informatics. Quantum computing with gates really starts when we go beyond a single qubit.

\section{Composite systems and the tensor product}\label{sec:tensor}
\subsection{Why a tensor product?}
\begin{origin}[: the tensor-product postulate]
If system $A$ can be in any of the states $\ket{i}_A$ ($i=1..d_A$) and system $B$ in any of the states $\ket\mu_B$ ($\mu=1..d_B$), then the
pair can be in any of the $d_Ad_B$ joint configurations $\ket i_A\ket\mu_B$, and, by the superposition principle, in any linear combination of
them. The state space is therefore of dimension $d_Ad_B$ (not $d_A+d_B$, as it would be for a classical description by a list of properties).
The tensor product $\Ha_A\otimes\Ha_B$ is precisely the space spanned by the products, with the scalar product
$(\bra i\bra\mu)(\ket j\ket\nu)=\braket ij\braket\mu\nu$. Probabilities of independent events multiply, which is consistent with
$|\braket{i\mu}{\psi_A\psi_B}|^2=|\braket i{\psi_A}|^2|\braket\mu{\psi_B}|^2$. For $n$ qubits the dimension is $2^n$.
\end{origin}

\subsection{Kronecker products of matrices and vectors}
Let $A$ be an $m\times n$ matrix and $B$ a $p\times q$ matrix. Their tensor (Kronecker) product $A\otimes B$ is the $mp\times nq$ matrix with elements
\begin{equation}
  (A\otimes B)_{p(r-1)+v,\ q(s-1)+w}=a_{rs}b_{vw},
\end{equation}
that is, the block matrix
\begin{equation}
  A\otimes B=\begin{pmatrix}a_{11}B&\cdots&a_{1n}B\\ \vdots&\ddots&\vdots\\ a_{m1}B&\cdots&a_{mn}B\end{pmatrix}.
\end{equation}
For example,
\begin{equation}
  \sigma_z\otimes\sigma_x=\begin{pmatrix}1\cdot\sigma_x&0\cdot\sigma_x\\0\cdot\sigma_x&-1\cdot\sigma_x\end{pmatrix}
  =\begin{pmatrix}0&1&0&0\\1&0&0&0\\0&0&0&-1\\0&0&-1&0\end{pmatrix}.
\end{equation}
The same rule applies to vectors, which fixes the ordering of the basis of the composite system:
\begin{equation}
  \ket0\otimes\ket0=\ket{00}=\begin{pmatrix}1\\0\end{pmatrix}\otimes\begin{pmatrix}1\\0\end{pmatrix}=\begin{pmatrix}1\\0\\0\\0\end{pmatrix},\quad
  \ket{01}=\begin{pmatrix}0\\1\\0\\0\end{pmatrix},\quad\ket{10}=\begin{pmatrix}0\\0\\1\\0\end{pmatrix},\quad\ket{11}=\begin{pmatrix}0\\0\\0\\1\end{pmatrix},
\end{equation}
and similarly $\ket{000}=(1,0,0,0,0,0,0,0)^T$. In general, the basis states of $n$ qubits are ordered as the binary numbers $00\dots0$, $00\dots1$, \dots, $11\dots1$
(big-endian convention, used throughout these notes). Useful rules:
\begin{equation}
  (A\otimes B)(C\otimes D)=AC\otimes BD,\quad(A\otimes B)^\dagger=A^\dagger\otimes B^\dagger,\quad\tr(A\otimes B)=\tr A\,\tr B,
\end{equation}
and the eigenvalues of $A\otimes B$ are the products $a_ib_j$ of the eigenvalues. An operator acting on qubit $k$ only, such as $\sigma^z_k$, means
$I\otimes\cdots\otimes\sigma^z\otimes\cdots\otimes I$ in the full $2^n$-dimensional space.

\section{Entanglement}
\begin{definition}
A pure state of $A$ and $B$ is a \emph{product} (separable) state if it can be written $\ket\psi=\ket{\psi_A}\otimes\ket{\psi_B}$. Otherwise it is \emph{entangled}.
For $N$ parts, a state is fully separable if it factorizes into a tensor product of $N$ single-part states.
\end{definition}
A state that can be factored is, for example,
\begin{equation}
  \Psi_{1+2}=\frac1{\sqrt2}\left(\ket0_1\ket1_2-\ket1_1\ket1_2\right)=\frac1{\sqrt2}\left(\ket0_1-\ket1_1\right)\ket1_2 .
\end{equation}
On the other hand, the following maximally entangled state cannot be factored:
\begin{equation}\label{eq:singlet}
  \Psi_{1+2}=\frac1{\sqrt2}\left(\ket0_1\ket1_2-\ket1_1\ket0_2\right).
\end{equation}
\begin{proof}
Suppose it can: $\frac1{\sqrt2}(\ket{01}-\ket{10})=(a\ket0+b\ket1)\otimes(c\ket0+d\ket1)=ac\ket{00}+ad\ket{01}+bc\ket{10}+bd\ket{11}$. Identifying coefficients,
\begin{equation}
  ac=0,\quad bd=0,\quad ad=\tfrac1{\sqrt2},\quad bc=-\tfrac1{\sqrt2}.
\end{equation}
The third and fourth equations require $a,b,c,d$ all non-zero, which contradicts the first two.
\end{proof}

\begin{proposition}[Entanglement criterion for two qubits]\label{prop:det}
The state $\ket\psi=\alpha\ket{00}+\beta\ket{01}+\gamma\ket{10}+\delta\ket{11}$ is a product state if and only if $\alpha\delta-\beta\gamma=0$.
\end{proposition}
\begin{proof}
Arrange the coefficients in the matrix $\Psi=\begin{pmatrix}\alpha&\beta\\ \gamma&\delta\end{pmatrix}$. A product state $(a,b)\otimes(c,d)$ has
$\Psi=\begin{pmatrix}a\\b\end{pmatrix}\begin{pmatrix}c&d\end{pmatrix}$, a matrix of rank one, hence $\det\Psi=0$. Conversely, a non-zero $2\times2$ matrix with
zero determinant has rank one, i.e.\ is an outer product of two vectors.
\end{proof}
The quantity $C=2|\alpha\delta-\beta\gamma|\in[0,1]$ is the \emph{concurrence} of the pure state; it is $0$ for product states and $1$ for maximally entangled
states such as \cref{eq:singlet}. The argument generalizes: a bipartite state is a product state iff its coefficient matrix has rank one. The rank is the Schmidt
number (\cref{app:schmidt}).

\paragraph{Consequence.} As soon as two subsystems 1 and 2 are entangled, there is no such thing as ``the wave function of subsystem 1'' alone. The whole is more
than the sum of its parts: all the information may be in the correlations, none in the parts. The description of a part requires the density matrix
(\cref{ch:density}). Entangled states are not rare: for random states of many qubits, product states form a set of measure zero.

\subsection{The Bell states}
The four maximally entangled states
\begin{equation}\label{eq:bellstates}
  \ket{\Phi^\pm}=\frac{\ket{00}\pm\ket{11}}{\sqrt2},\qquad\ket{\Psi^\pm}=\frac{\ket{01}\pm\ket{10}}{\sqrt2}
\end{equation}
form an orthonormal basis of two qubits, the \emph{Bell basis}. $\ket{\Psi^-}$ is the spin \emph{singlet} (total spin $0$); the other three form the
\emph{triplet} (total spin $1$). They are all related by local operations on one qubit: $\ket{\Phi^-}=(Z\otimes I)\ket{\Phi^+}$,
$\ket{\Psi^+}=(X\otimes I)\ket{\Phi^+}$, $\ket{\Psi^-}=(XZ\otimes I)\ket{\Phi^+}$ (up to sign). They are created from product states by the circuit
\begin{center}
\begin{quantikz}
\lstick{$\ket{x}$} & \gate{H} & \ctrl{1} & \qw \\
\lstick{$\ket{y}$} & \qw & \targ{} & \qw
\end{quantikz}
\qquad$\ket{00}\to\ket{\Phi^+},\ \ket{01}\to\ket{\Psi^+},\ \ket{10}\to\ket{\Phi^-},\ \ket{11}\to\ket{\Psi^-}$.
\end{center}
For example, $H\otimes I$ maps $\ket{00}$ to $(\ket{00}+\ket{10})/\sqrt2$, and the CNOT (\cref{sec:cnot}) flips the second qubit in the second term, giving
$(\ket{00}+\ket{11})/\sqrt2$. Running the circuit backwards (CNOT, then $H$) maps the Bell basis onto the computational basis: this is how a \emph{Bell measurement}
is performed on a quantum computer.

\subsection{Multipartite entanglement\starred}
For three qubits, two inequivalent kinds of genuine tripartite entanglement exist:
\begin{equation}
  \ket{\mathrm{GHZ}}=\frac{\ket{000}+\ket{111}}{\sqrt2},\qquad\ket{W}=\frac{\ket{001}+\ket{010}+\ket{100}}{\sqrt3}.
\end{equation}
Tracing out one qubit of GHZ leaves a separable (classically correlated) state, while for W the remaining two qubits are still entangled. No local operations,
even assisted by classical communication and with a finite probability of success, can convert one into the other (Dür, Vidal, Cirac 2000).

\section{Quantum parallelism and entanglement}\label{sec:parallelism}
Why is entanglement important for quantum computing? Because it makes the link between the input and the output.

The idea of \emph{quantum parallelism} is based on the possibility for a qubit to be in a superposition of two states. It is therefore possible to perform a
computation for all possible inputs (input qubits in their superpositions) and obtain all possible outputs in a single run. Given a classical function
$f:\{0,1\}^n\to\{0,1\}^m$, the reversible implementation $U_f\ket x\ket y=\ket x\ket{y\oplus f(x)}$ (\cref{sec:reversible}) gives
\begin{equation}
  U_f\left(\frac1{\sqrt{2^n}}\sum_x\ket x\right)\ket0=\frac1{\sqrt{2^n}}\sum_x\ket x\ket{f(x)}.
\end{equation}
A measurement at the end gives one single result, but quantum algorithms enhance the probability to measure the result ``of interest''. For example, one can look
for the solutions of an equation, the output being its right-hand side. All possible inputs give all possible values, but interference can increase the probability
to measure the correct output. How do we then know to which input it corresponds, i.e.\ to which values of the unknowns? Only thanks to entanglement: each output
is entangled with the corresponding inputs, so that when we observe a value of the output qubits, the input qubits collapse onto the states corresponding to that output.

\begin{example}[Multiplication of two bits]
We want
\begin{equation}
  \ket{00}\to\ket0,\quad\ket{01}\to\ket0,\quad\ket{10}\to\ket0,\quad\ket{11}\to\ket1 .
\end{equation}
The state of the system should be a superposition of the four 3-qubit states $\ket{000}$, $\ket{010}$, $\ket{100}$, $\ket{111}$ (2 input qubits and 1 output qubit)
and no other terms. This state is entangled: a product of single-qubit states containing $\ket{000}$ and $\ket{111}$ would also contain terms such as $\ket{001}$ or
$\ket{110}$. Measuring the output and finding $1$ projects the inputs onto $\ket{11}$; finding $0$ projects them onto $(\ket{00}+\ket{01}+\ket{10})/\sqrt3$.
The circuit is a Toffoli gate (\cref{sec:toffoli}) with the output qubit as target.
\end{example}

\begin{remark}
Entanglement is necessary but not sufficient for a quantum speed-up. By the Gottesman--Knill theorem, circuits made only of $H$, $S$ and CNOT gates can
create highly entangled states but are efficiently simulable classically. Conversely, a pure-state computation whose entanglement stays bounded (small Schmidt rank
across every cut) can be simulated efficiently with tensor networks (Vidal 2003). Since entanglement is crucial, there is a large effort to create and maximize it,
and this requires being able to measure it and describe its decay by decoherence (\cref{ch:density,ch:open}).
\end{remark}

\section{Exercises}
\begin{exercise}
Which of the following states are entangled? $\frac12(\ket{00}+\ket{01}+\ket{10}+\ket{11})$; $\frac12(\ket{00}+\ket{01}+\ket{10}-\ket{11})$;
$\frac1{\sqrt2}(\ket{00}+\ii\ket{11})$. Compute their concurrence.
\end{exercise}
\begin{exercise}
Show that $\ket{\Psi^-}$ has the same form in every basis: $\ket{\Psi^-}=\frac1{\sqrt2}(\ket{\vec n,+}\ket{\vec n,-}-\ket{\vec n,-}\ket{\vec n,+})$ for any axis $\vec n$.
(Hint: $(U\otimes U)\ket{\Psi^-}=\det U\,\ket{\Psi^-}$ for any $2\times2$ matrix $U$.)
\end{exercise}
